% ==============================================================================
% Course Syllabus: MSQF 535 Global Financial Management
% Formatted to fit cleanly on a single official syllabus page
% ==============================================================================

\newpage
\thispagestyle{plain}

\begin{center}
    {\large\textbf{MSQF 535: GLOBAL FINANCIAL MANAGEMENT}} \hfill {\large\textbf{CREDITS: 4}}
\end{center}
\vspace{-0.7em}
\hrule height 0.8pt
\vspace{0.35em}

\noindent\textbf{Course Objectives:}
\begin{itemize}[leftmargin=1.2em, label=$\bullet$, noitemsep, topsep=1pt, parsep=0pt]
    \item \small\textit{The aim of this course is to provide an understanding of the fundamental concepts and managerial issues pertaining to international finance.}\normalsize
\end{itemize}

\vspace{0.25em}
\noindent\textbf{Course Outcomes (COs):}
\begin{description}[leftmargin=1.1cm, labelsep=0.3em, itemsep=0.5pt, topsep=1pt, parsep=0pt, font=\normalfont\small\bfseries]
    \small
    \item[CO1:] \textbf{Understand the global financial landscape:} Comprehend the roles and functions of international financial institutions and development banks like the World Bank, IMF, and ADB, as well as the process of internationalization.
    \item[CO2:] \textbf{Analyze foreign exchange markets and rates:} Explain the mechanisms of the foreign exchange market, including spot and forward markets, cross-rates, bid-ask spreads, and apply theories determining exchange rates.
    \item[CO3:] \textbf{Manage foreign exchange exposure:} Evaluate and manage various types of foreign exchange exposure (translation, transaction, and economic) using different methods and strategies.
    \item[CO4:] \textbf{Assess multinational firm's financial management:} Analyze foreign direct investment decisions, cost of capital, capital structure, budgeting, cash management, country risk, and international taxation for multinational firms.
    \item[CO5:] \textbf{Evaluate international financing operations:} Understand the characteristics and instruments of Eurocurrency markets, interest rate and currency swaps, depository receipts (GDR and ADR), and the implications of the Euro for India.
    \normalsize
\end{description}

\vspace{0.2em}
\noindent\textbf{Unit-I: International Finance Overview \& Institutions}\\
\footnotesize
International Finance --- overview --- International Financial Institutions/Development Banks --- World Bank --- IBRD --- IDA --- IFC --- MIGA --- International Monetary Fund --- Special Drawing Rights --- Asian Development Bank --- Internationalization process.
\normalsize

\vspace{0.2em}
\noindent\textbf{Unit-II: The Foreign Exchange Market}\\
\footnotesize
The Foreign Exchange Market --- SWIFT --- Arbitrage --- Spot market --- Forward market --- Cross rates of exchange --- Bid --- Ask spreads --- Balance of payments --- Foreign exchange rates --- Theories of Foreign Exchange Rate.
\normalsize

\vspace{0.2em}
\noindent\textbf{Unit-III: Foreign Exchange Exposure and Management}\\
\footnotesize
Foreign exchange exposure and management --- Management of translation exposure --- Methods --- Management of transaction exposure --- Management of economic exposure --- Methods --- Strategies.
\normalsize

\vspace{0.2em}
\noindent\textbf{Unit-IV: Financial Management of the Multinational Firm}\\
\footnotesize
Financial Management of the Multinational Firm --- Foreign direct investment --- Cost of capital and capital structure of the multinational firm --- Multinational capital budgeting --- Multinational cash management --- Country Risk Analysis --- International Taxation.
\normalsize

\vspace{0.2em}
\noindent\textbf{Unit-V: Financing Foreign Operations}\\
\footnotesize
Financing Foreign Operations --- Eurocurrency markets --- Instruments --- Interest rate swaps --- Currency swaps and its pricing --- Depository receipts --- GDR and ADR --- Euro and its implications for India.
\normalsize

\vspace{0.25em}
\noindent\textbf{Books for References:}
\begin{enumerate}[leftmargin=1.2em, label=\arabic*., itemsep=0pt, topsep=1pt, parsep=0pt]
    \scriptsize
    \item Buckley, A. (1992): \textit{Multinational Finance}, Prentice Hall, New Jersey.
    \item Levich, R. M. (2001): \textit{International Financial Markets: Prices and Policies}, McGraw Hill, New York.
    \item Vij, M. (2001): \textit{Multinational Financial Management}, Excel Books, New Delhi.
    \item Shapiro, A. C. (1996): \textit{Multinational Financial Management}, Prentice Hall of India, New Delhi.
    \item Apte, P. G. (1999): \textit{International Financial Management}, Tata McGraw Hill Publishing Company Ltd.
    \item Jain, P. K., et al. (1998): \textit{International Financial Management}, Macmillan, New Delhi.
    \item Eun, C. S., \& Resnick, B. G. (2001): \textit{International Financial Management}, Irwin McGraw Hill, Singapore.
\end{enumerate}

\vspace{0.2em}
\begin{center}
\scriptsize
\setlength{\tabcolsep}{6pt}
\renewcommand{\arraystretch}{0.95}
\begin{tabular}{|c|c|c|c|c|c|}
\hline
\textbf{COs} & \textbf{PO1} & \textbf{PO2} & \textbf{PO3} & \textbf{PO4} & \textbf{PO5} \\
\hline
\textbf{CO1} & XXX & X & X & X & XX \\
\hline
\textbf{CO2} & XX & X & XX & X & X \\
\hline
\textbf{CO3} & X & XX & XX & X & X \\
\hline
\textbf{CO4} & XX & X & XX & XX & X \\
\hline
\textbf{CO5} & XXX & X & X & XX & XX \\
\hline
\end{tabular}
\end{center}
\normalsize

\newpage
